\documentclass{amsart}
% For dealing with references we use the comment environment
\usepackage{verbatim}
\newenvironment{reference}{\comment}{\endcomment}
%\newenvironment{reference}{}{}
\newenvironment{slogan}{\comment}{\endcomment}
\newenvironment{history}{\comment}{\endcomment}

% For commutative diagrams we use Xy-pic
\usepackage[all]{xy}

% We use 2cell for 2-commutative diagrams.
\xyoption{2cell}
\UseAllTwocells

% We use multicol for the list of chapters between chapters
\usepackage{multicol}

% This is generally recommended for better output
\usepackage{lmodern}
\usepackage[T1]{fontenc}

% For cross-file-references

% Package for hypertext links:
\usepackage{hyperref}

% For any local file, say "hello.tex" you want to link to please

% Theorem environments.
%
\theoremstyle{plain}
\newtheorem{theorem}[subsection]{Theorem}
\newtheorem{proposition}[subsection]{Proposition}
\newtheorem{lemma}[subsection]{Lemma}

\theoremstyle{definition}
\newtheorem{definition}[subsection]{Definition}
\newtheorem{example}[subsection]{Example}
\newtheorem{exercise}[subsection]{Exercise}
\newtheorem{situation}[subsection]{Situation}

\theoremstyle{remark}
\newtheorem{remark}[subsection]{Remark}
\newtheorem{remarks}[subsection]{Remarks}

\numberwithin{equation}{subsection}

% Macros
%
\def\lim{\mathop{\mathrm{lim}}\nolimits}
\def\colim{\mathop{\mathrm{colim}}\nolimits}
\def\Spec{\mathop{\mathrm{Spec}}}
\def\Hom{\mathop{\mathrm{Hom}}\nolimits}
\def\Ext{\mathop{\mathrm{Ext}}\nolimits}
\def\SheafHom{\mathop{\mathcal{H}\!\mathit{om}}\nolimits}
\def\SheafExt{\mathop{\mathcal{E}\!\mathit{xt}}\nolimits}
\def\Sch{\mathit{Sch}}
\def\Mor{\mathop{\mathrm{Mor}}\nolimits}
\def\Ob{\mathop{\mathrm{Ob}}\nolimits}
\def\Sh{\mathop{\mathit{Sh}}\nolimits}
\def\NL{\mathop{N\!L}\nolimits}
\def\CH{\mathop{\mathrm{CH}}\nolimits}
\def\proetale{{pro\text{-}\acute{e}tale}}
\def\etale{{\acute{e}tale}}
\def\QCoh{\mathit{QCoh}}
\def\Ker{\mathop{\mathrm{Ker}}}
\def\Im{\mathop{\mathrm{Im}}}
\def\Coker{\mathop{\mathrm{Coker}}}
\def\Coim{\mathop{\mathrm{Coim}}}

% Boxtimes
%
\DeclareMathSymbol{\boxtimes}{\mathbin}{AMSa}{"02}

%
% Macros for moduli stacks/spaces
%
\def\QCohstack{\mathcal{QC}\!\mathit{oh}}
\def\Cohstack{\mathcal{C}\!\mathit{oh}}
\def\Spacesstack{\mathcal{S}\!\mathit{paces}}
\def\Quotfunctor{\mathrm{Quot}}
\def\Hilbfunctor{\mathrm{Hilb}}
\def\Curvesstack{\mathcal{C}\!\mathit{urves}}
\def\Polarizedstack{\mathcal{P}\!\mathit{olarized}}
\def\Complexesstack{\mathcal{C}\!\mathit{omplexes}}
% \Pic is the operator that assigns to X its picard group, usage \Pic(X)
% \Picardstack_{X/B} denotes the Picard stack of X over B
% \Picardfunctor_{X/B} denotes the Picard functor of X over B
\def\Pic{\mathop{\mathrm{Pic}}\nolimits}
\def\Picardstack{\mathcal{P}\!\mathit{ic}}
\def\Picardfunctor{\mathrm{Pic}}
\def\Deformationcategory{\mathcal{D}\!\mathit{ef}}

\setlength{\emergencystretch}{3em}
\begin{document}
\title[Stacks Project, Tag 0EZK]{271 --- Descent of etale sheaf sections along submersive morphisms with connected geometric fibres}
\maketitle
\noindent Copyright (C) 2005--2025 Johan de Jong. Stacks Project authors. Source: \href{https://stacks.math.columbia.edu/tag/0EZK}{Tag 0EZK}.

\noindent Editorial difficulty note (not author text). Difficulty H2: One must turn equality of stalk values into equality on entire geometric fibres, using connectedness of the pulled-back constant sheaf. The saturated maximal equality locus must then descend to an open subset through submersiveness, allowing local sections to glue globally..

\noindent Editorial provenance note (not author text). History: excerpt from revision \texttt{a0444\allowbreak 6e57e\allowbreak c1fbc\allowbreak 252a8\allowbreak 71afc\allowbreak ec775\allowbreak 2fb28\allowbreak 07b14}, etale-cohomology.tex, lines 5076--5140. Reference presentation was adapted; original-author context and core support blocks were retained or added. The mathematical wording is unchanged. Licensed under GNU FDL 1.2 or later; no invariant sections or cover texts. See ../licenses/COPYING. Context blocks reproduce author definitions and supporting results; they are not additional counted problems.

\begin{lemma}
\label{lemma-sections-upstairs-submersive}
Let $S$ be a scheme. Let $f : X \to S$ be a morphism such that
\begin{enumerate}
\item $f$ is submersive, and
\item the geometric fibres of $f$ are connected.
\end{enumerate}
Let $\mathcal{F}$ be a sheaf on $S_\etale$.
Then $\Gamma(S, \mathcal{F}) = \Gamma(X, f^{-1}_{small}\mathcal{F})$.
\end{lemma}

\begin{proof}
There is a canonical map
$\Gamma(S, \mathcal{F}) \to \Gamma(X, f_{small}^{-1}\mathcal{F})$.
Since $f$ is surjective (because its fibres are connected) we see that
this map is injective.

\medskip\noindent
To show that the map is surjective, let
$\tau \in \Gamma(X, f_{small}^{-1}\mathcal{F})$.
It suffices to find
an \'etale covering $\{U_i \to S\}$ and
sections $\sigma_i \in \mathcal{F}(U_i)$
such that $\sigma_i$ pulls back to $\tau|_{X \times_S U_i}$.
Namely, the injectivity shown above guarantees
that $\sigma_i$ and $\sigma_j$ restrict to the same
section of $\mathcal{F}$ over $U_i \times_S U_j$.
Thus we obtain a unique section $\sigma \in \mathcal{F}(S)$
which restricts to $\sigma_i$ over $U_i$.
Then the pullback of $\sigma$ to $X$ is $\tau$
because this is true locally.

\medskip\noindent
Let $\overline{x}$ be a geometric point of $X$ with image $\overline{s}$
in $S$. Consider the image of $\tau$ in the stalk
$$
(f_{small}^{-1}\mathcal{F})_{\overline{x}} = \mathcal{F}_{\overline{s}}
$$
See Lemma \href{https://stacks.math.columbia.edu/tag/03Q1}{lemma stalk pullback (Tag 03Q1)}.
We can find an \'etale neighbourhood $U \to S$ of $\overline{s}$
and a section $\sigma \in \mathcal{F}(U)$ mapping to this image
in the stalk. Thus after replacing $S$ by $U$ and $X$ by $X \times_S U$
we may assume there exits a section $\sigma$ of $\mathcal{F}$ over $S$
whose image in $(f_{small}^{-1}\mathcal{F})_{\overline{x}}$ is the
same as $\tau$.

\medskip\noindent
By Lemma \href{https://stacks.math.columbia.edu/tag/09XM}{lemma where sections are equal (Tag 09XM)}
there exists a maximal open $W \subset X$ such that
$f_{small}^{-1}\sigma$ and $\tau$ agree over $W$
and the formation of $W$ commutes with further pullback.
Observe that the pullback of $\mathcal{F}$ to the
geometric fibre $X_{\overline{s}}$ is the pullback
of $\mathcal{F}_{\overline{s}}$ viewed as a sheaf on
$\overline{s}$ by $X_{\overline{s}} \to \overline{s}$.
Hence we see that $\tau$ and $\sigma$ give sections
of the constant sheaf with value $\mathcal{F}_{\overline{s}}$
on $X_{\overline{s}}$ which agree in one point. Since
$X_{\overline{s}}$ is connected by assumption, we conclude
that $W$ contains $X_s$. The same argument for different
geometric fibres shows that $W$ contains every fibre it meets.
Since $f$ is submersive, we conclude that $W$ is the inverse
image of an open neighbourhood of $s$ in $S$.
This finishes the proof.
\end{proof}
\end{document}
